\documentclass[english]{article}
\usepackage[T1]{fontenc}
\usepackage[utf8]{inputenc}
\setcounter{secnumdepth}{-2}
\setlength{\parindent}{0bp}
\usepackage{amsmath}
\usepackage{amsthm}
\usepackage{amssymb}

\makeatletter
\theoremstyle{plain}
\newtheorem{thm}{\protect\theoremname}
\theoremstyle{definition}
\newtheorem{problem}[thm]{\protect\problemname}
\theoremstyle{definition}
\newtheorem*{sol*}{\protect\solutionname}

\@ifundefined{date}{}{\date{}}
\usepackage{graphicx}
\usepackage{geometry}
\newcommand{\limi}{\underset{n\rightarrow\infty}{\lim}}
\newcommand{\limxz}{\underset{x\rightarrow x_{0}}{\lim}}
\newcommand{\limfxz}{\underset{x\rightarrow x_{0}}{\lim}f(x)}
\newcommand{\limfxm}{\underset{x\rightarrow x_{0}^{-}}{\lim}f(x)}
\newcommand{\limfxp}{\underset{x\rightarrow x_{0}^{+}}{\lim}f(x)}
\newcommand{\limxm}{\underset{x\rightarrow x_{0}^{-}}{\lim}}
\newcommand{\limxp}{\underset{x\rightarrow x_{0}^{+}}{\lim}}
\newcommand{\sqrm}{M_{m\times m}(\mathbb{F})}
\newcommand{\seqa}{(a_{n})_{n=1}^{\infty}}
\newcommand{\seqx}{(x_{n})_{n=1}^{\infty}}
\newcommand{\funcdr}{f:D\rightarrow\mathbb{R}}
\newcommand{\der}{\limxz\frac{f(x)-f(x_{0})}{x-x_{0}}}
\usepackage[all]{pdfprivacy}

\makeatother

\usepackage{babel}
\providecommand{\problemname}{Problem}
\providecommand{\solutionname}{Solution}
\providecommand{\theoremname}{Theorem}

\begin{document}
\newgeometry{left=2cm, right=2cm}
\title{{\Large\vspace{-6em}Semester 2026B, Quiz 1\vspace{-3em}}}
\maketitle
\begin{abstract}
Quiz 1 in the course Infinitesimal Calculus (2), semester 2026 B.
\\

This is a quiz that I translated and solved.

The original quiz was written by the Infi Calculus 2 course lecturers
at HUJI.
\end{abstract}
\medskip{}

\rule[0.5ex]{1\columnwidth}{1pt}

\medskip{}

\begin{problem}
Find a $q\in\mathbb{Q}\text{ s.t.}\left|e^{-1}-q\right|<10^{-2}$.
\end{problem}

\begin{sol*}
Our goal is to find a rational approximation for the function $e^{x}$
evaluated at $x=-1$, with an error strictly less than $10^{-2}$.

The closest point to $x=-1$ which we can easily evaluate is $x=0$
since $e^{0}=1$.

The Maclaurin polynomial (the Taylor polynomial at 0) of degree $n$
of $e^{x}$ is 
\[
P_{n,e^{x},0}(x)=\underset{k=0}{\overset{n}{\sum}}\frac{f^{(k)}(0)}{k!}x^{k}=e^{0}+e^{0}\cdot x+\frac{e^{0}}{2!}\cdot x^{2}+\dots+\frac{e^{0}}{n!}\cdot x^{n}=1+x+\frac{x^{2}}{2!}+\dots+\frac{x^{n}}{n!}
\]

Let $n\in\mathbb{N}$. To apply Taylor's Theorem on $e^{x}$ in the
interval $\left[-1,0\right]$, $e^{x}$ must have a continuous $n$-th
derivative on $\left[-1,0\right]$ and be differentiable $n+1$ times
on $\left(-1,0\right)$.

$e^{x}$ is differentiable $n$ times on $\left[-1,0\right]$ with
its $n$-th derivative being $e^{x}$, which is continuous on $\mathbb{R}$
and in particular on $\left[-1,0\right]$.

$e^{x}$ is also differentiable $n+1$ times on $\left(-1,0\right)$.

In fact, $e^{x}$ is a smooth function (it is infinitely differentiable)
on $\mathbb{R}$.

Hence, the Taylor Theorem with the Lagrange Remainder applies and
there exists a point $c\in\left(-1,0\right)$ such that 
\[
R_{n,e^{x},0}(-1)=e^{-1}-P_{n,e^{x},0}(-1)=\frac{e^{c}}{\left(n+1\right)!}\left(-1\right)^{n+1}
\]
 This means that 
\[
\left|e^{-1}-P_{n,e^{x},0}(-1)\right|\le\left|\frac{e^{c}}{\left(n+1\right)!}\left(-1\right)^{n+1}\right|=\left|\frac{e^{c}}{\left(n+1\right)!}\right|\underset{\left(*\right)}{\le}\frac{1}{\left(n+1\right)!}
\]

$\left(*\right)$ since $e^{x}$ is a monotonically increasing function,
$\forall x\in\left(-1,0\right),e^{x}\le e^{0}=1$.

To guarantee our error is strictly less than $10^{-2}$, we need $\frac{1}{(n+1)!}<\frac{1}{100}$,
which implies $(n+1)!>100$.

Since $4!=24$ and $5!=120$, the minimum degree required is $n=4$.

Therefore, setting the degree $n=4$ gives the rational approximation:
\[
P_{4,e^{x},0}\left(-1\right)=1-1+\frac{1}{2}-\frac{1}{6}+\frac{1}{24}=\frac{3}{8}\in\mathbb{Q}
\]

which satisfies the error bound:
\[
\left|e^{-1}-P_{4,e^{x},0}(-1)\right|\le\frac{1}{5!}<\frac{1}{100}=10^{-2}
\]

as required.
\end{sol*}
\medskip{}

\rule[0.5ex]{1\columnwidth}{1pt}

\medskip{}

\newpage{}
\begin{problem}
Find all antiderivatives of $f(x)=\frac{x^{2}-4x}{x^{2}-3x+2}$ on
$\left(1,2\right)$.
\end{problem}

\begin{sol*}
Since the degree of the numerator is greater than or equal to the
degree of the denominator, our first step is to perform polynomial
division with a remainder.

We can achieve this by expressing the numerator in terms of the denominator.
\[
f(x)=\frac{x^{2}-3x+2-x-2}{x^{2}-3x+2}=1+\frac{-x-2}{x^{2}-3x+2}
\]

Next, we focus on finding the antiderivatives of the remainder, $g(x)=\frac{-x-2}{x^{2}-3x+2}$,
on the interval $(1,2)$.

As a best practice, we first check if the numerator is a scalar multiple
of the denominator's derivative, which would allow for direct logarithmic
integration.
\[
\left(x^{2}-3x+2\right)'=2x-3
\]

Since this is not the case, we proceed with partial fraction decomposition. 

We first split the denominator into coprime polynomials. 
\[
x^{2}-3x+2=\left(x-2\right)\left(x-1\right)
\]

Hence
\[
g(x)=\frac{-x-2}{\left(x-2\right)\left(x-1\right)}
\]
\[
\frac{-x-2}{\left(x-2\right)\left(x-1\right)}=\frac{A}{x-2}+\frac{B}{x-1}
\]
\[
-x-2=\left(x-1\right)A+\left(x-2\right)B
\]
\[
\left(A+B\right)x-A-2B=-x-2
\]

And since two polynomials are equal if and only if their corresponding
coefficients are equal,
\[
\begin{cases}
A+B=-1\\
-A-2B=-2
\end{cases}\Rightarrow\begin{cases}
A=-4\\
B=3
\end{cases}
\]
\[
g(x)=\frac{-4}{x-2}+\frac{3}{x-1}
\]
Now we can return to our original function, $f(x)=1+g(x)$, and find
the antiderivatives term by term.

The antiderivative of 1 is $F_{1}(x)=x+C$. 

The antiderivative of $\frac{-4}{x-2}$ is 
\[
F_{2}(x)=-4\ln\left|x-2\right|+C
\]
 and the antiderivative of $\frac{3}{x-1}$ is 
\[
F_{3}(x)=3\ln\left|x-1\right|+C
\]

Therefore, the general antiderivative of the function is 
\[
F(x)=x-4\ln\left|x-2\right|+3\ln\left|x-1\right|+C
\]

For $x\in(1,2)$, we have $x-2<0$ and $x-1>0$. Thus, the antiderivative
simplifies to:
\[
F(x)=x-4\ln\left(2-x\right)+3\ln\left(x-1\right)+C
\]
\end{sol*}

\end{document}
